\section{From uniformly generated cone ideals to the original envelope}

The following argument is the analytic assembly step in the multiplier-ideal
proof of the envelope property.  It uses the actual analytification of a
projective subvariety, the obstacle ideals of
\cref{def:sg-Jm}, and the original Fubini--Study system.  The algebraic
hypotheses below must be established for the multiplier ideals under
consideration; the criterion alone makes no assertion that an arbitrary
smooth projective variety supplies those ideals.

\begin{definition}[Bounded generators on the punctured cone]
\label{def:c0n-bounded-cone-generators}
\uses{conv:sg-setting,lem:sg-phi-basic}
\lean{EnvelopeProperty.ProjSubvariety.HasBoundedConeGenerators}
\leanok
Let $X\subset\mathbf P^M_k$ be an integral projective subvariety, let
$T=(k[x_0,\ldots,x_M]/I_X)[t]$, and let $\mathfrak m_+$ be its irrelevant
ideal.  An ideal $B\subset T$ has bounded cone generators of degree $e$ if
there is a set $G\subset T$ such that
\[
 (G)\subset B\subset (G)^{\mathrm{sat}}
\]
and every $g\in G$ has the form $t^\ell\bar s$, where $s$ is homogeneous of
some degree $b\le e$.  Saturation is with respect to $\mathfrak m_+$.
\end{definition}

\begin{lemma}[Fubini--Study functions from bounded cone generators]
\label{lem:c0n-bounded-cone-fs}
\uses{def:c0n-bounded-cone-generators,lem:sg-phi-basic}
\lean{EnvelopeProperty.ProjSubvariety.isFS_coneIdealApprox_of_bounded_generators}
\leanok
Let $B\ne0$ have bounded cone generators of degree $md+c$, where $m>0$.
Then
\[
 u_m=m^{-1}\phi_B\in\mathcal H(X,d+c/m).
\]
\end{lemma}
\begin{proof}
Lift each generator $t^\ell\bar s$ to $t^\ell s$ in
$k[x_0,\ldots,x_M,t]$.  Let $B'$ be the ideal generated by the lifts.
Saturation does not change an ideal potential on $X^{\mathrm{an}}$, and
neither does adding $I_X[t]$.  Hence $\phi_B=\phi_{B'}$ on
$X^{\mathrm{an}}$.  At least one lifted generator has nonzero image in $T$:
otherwise $(G)=0$, whereas the saturation of zero is zero because $T$ is a
domain and $\mathfrak m_+$ contains a nonzero coordinate.  The coefficients
of the lifts are homogeneous of degree at most $md+c$.  The finite-generator
formula for ideal potentials, padding homogeneous forms with coordinates,
therefore gives $\phi_{B'}\in\mathcal H(X,md+c)$.  Multiplication by $m^{-1}$
gives the asserted class because $(md+c)/m=d+c/m$.
\end{proof}

\begin{lemma}[From element orders to cone potentials]
\label{lem:c0n-orders-to-potentials}
\uses{lem:sg-phi-basic,def:sg-div}
\lean{EnvelopeProperty.ProjSubvariety.coneIdealPotential_le_of_element_orders}
\leanok
Let a divisorial point be represented by a discrete valuation ring $V$ and
a positive rational number $a$, so that $\log|g|=-a\operatorname{ord}_V(\bar g)$.
Let $J\ne0$ and $B$ be cone ideals, let $\lambda\ge0$, and let $A$ be real.
Suppose that whenever $o\le\operatorname{ord}_V(g)$ for every $g\in J$,
every nonzero $f\in B$ satisfies
$\lambda o-A\le\operatorname{ord}_V(f)$.  Then
\[
 \phi_B(v)\le\lambda\phi_J(v)+aA.
\]
\end{lemma}
\begin{proof}
Since $J\ne0$, the set of finite orders of its nonzero elements is a
nonempty subset of $\mathbf N$.  Let $o$ be its least element and choose
$g_0\in J$ of order $o$.  Then $-ao\le\phi_J(v)$.  If $f\in B$ is
nonzero and has order $\nu$, the hypothesis gives
$-a\nu\le-a\lambda o+aA\le\lambda\phi_J(v)+aA$.
The same bound holds for the zero element, whose logarithmic value is
$-\infty$.  Lift all elements of $B$ to the polynomial ring and take the
supremum of their logarithmic values.  This supremum is $\phi_B(v)$,
which proves the assertion.
\end{proof}

\begin{lemma}[Normalized cone ideal bounds]
\label{lem:c0n-cone-ideal-bounds}
\source{bj22-trivval:thm:contenvlisse}
\uses{prop:sg-envelope-Jm,lem:sg-phi-basic}
\lean{EnvelopeProperty.ProjSubvariety.coneIdealApprox_mul_le,EnvelopeProperty.ProjSubvariety.envelope_le_coneIdealApprox,EnvelopeProperty.ProjSubvariety.exists_coneIdealApprox_le_envelope_add}
\leanok
Let $d>0$, let $\phi$ be lower semicontinuous and bounded below with
$\phi(\star)\le0$, and write $J_q=J_q(\phi)T$.  Suppose that ideals $B_m$
satisfy $J_m\subset B_m$ for $m>0$ and
$B_{ml}\subset(B_m^l)^{\mathrm{sat}}$ for $m,l>0$.
Then $u_{ml}\le u_m$ and $P_d(\phi)\le u_m$.
Suppose moreover that for every divisorial $v$ there is $A_v\ge0$ such
that, for each $m>0$, some $q>0$ satisfies
\[
 \phi_{B_m}(v)\le (m/q)\phi_{J_q}(v)+A_v.
\]
Then $u_m(v)\le P_d(\phi)(v)+A_v/m$.
\end{lemma}
\begin{proof}
The power and saturation identities for ideal potentials yield
$\phi_{B_{ml}}\le l\phi_{B_m}$.  Division by $ml$ gives $u_{ml}\le u_m$.
For every $q>0$, gradedness of the obstacle ideals and the inclusion into
$B_{qm}$ give
\[
 q^{-1}\phi_{J_q}\le(qm)^{-1}\phi_{J_{qm}}
 \le u_{qm}\le u_m.
\]
Taking the supremum over $q$ and using the obstacle-ideal formula for the
envelope proves $P_d(\phi)\le u_m$.  Finally divide the displayed
valuative estimate by $m$.  The normalized obstacle potential
$q^{-1}\phi_{J_q}$ is at most the envelope, giving the final inequality.
\end{proof}

\begin{theorem}[Envelope assembly from the cone ideal estimates]
\label{thm:c0n-cone-envelope-assembly}
\source{bj22-trivval:thm:contenvlisse}
\uses{lem:c0n-bounded-cone-fs,lem:c0n-cone-ideal-bounds,thm:pg-psh2,thm:sg-comparison}
\lean{EnvelopeProperty.ProjSubvariety.isPSH_envelope_of_coneIdealEstimates}
\leanok
Assume that $k$ is algebraically closed.  Let $d>0$ and let $\phi\le0$ be
piecewise linear on $X^{\mathrm{an}}$.  Suppose that $B_m$ satisfy all the
hypotheses of \cref{lem:c0n-cone-ideal-bounds} and that, for one fixed
$c\in\mathbf N$, $B_m$ has bounded cone generators of degree $md+c$ for
every $m>0$.  Then $P_d(\phi)$ is $d$-plurisubharmonic.
\end{theorem}
\begin{proof}
Put $U=\inf_j u_{j!}$.  The sequence $u_{j!}$ decreases, each term is
plurisubharmonic in class $d+c/j!$, and these classes tend to $d$.
The continuous function $\phi$ is bounded below, say by $C$.
The constant $C$ is a competitor for $P_d(\phi)$, so
$U\ge P_d(\phi)\ge C$.  Thus $U$ is not identically $-\infty$.
The decreasing-limit theorem proves that $U$ is $d$-plurisubharmonic.
At a divisorial point $v$ we have
\[
 U(v)\le u_{j!}(v)\le P_d(\phi)(v)+A_v/j!
\]
for every $j$.  Given $\varepsilon>0$, choose $j>A_v/\varepsilon$;
since $j!\ge j$, the error is at most $\varepsilon$.  Letting
$\varepsilon$ decrease to zero gives $U(v)\le P_d(\phi)(v)\le\phi(v)$.
Divisorial comparison with the continuous obstacle yields $U\le\phi$
everywhere.  Therefore $U$ is a competitor and $U\le P_d(\phi)$.
The opposite inequality was already proved, so $U=P_d(\phi)$.
\end{proof}

\begin{remark}
The arithmetic and decreasing-limit argument follows the multiplier-ideal
comparison in the proof of \cite[Theorem~5.20]{BJ22}.  The route here uses
obstacle ideals on the homogeneous cone and divisorial comparison rather
than a uniform estimate on a resolved test configuration.  Its geometric
input is precisely the production of the ideals and their bounded
generators and discrepancy estimates; this input is not an envelope
conclusion.
\end{remark}
